% !TEX program = xelatex
\documentclass[11pt]{ctexart}
\usepackage[a4paper,left=2.0cm,right=2.0cm,top=2.0cm,bottom=2.2cm]{geometry}
\usepackage{amsmath,amssymb}
\usepackage{xcolor}
\usepackage{array}
\usepackage{enumitem}

\setlist{nosep,leftmargin=1.6em}
\linespread{1.23}
\setlength{\parskip}{3pt}
\setlength{\parindent}{0pt}

\definecolor{deepblue}{RGB}{12,82,160}
\definecolor{deepnavy}{RGB}{18,60,104}
\definecolor{band}{RGB}{242,246,251}
\definecolor{rulegray}{RGB}{200,212,228}

\newcommand{\hdr}[1]{%
  \par\vspace{10pt}%
  \noindent{\color{deepblue}\large\bfseries #1}\par\vspace{2pt}%
  \noindent{\color{deepblue}\rule{\linewidth}{1.1pt}}\par\vspace{6pt}}

\newcommand{\resultband}[1]{%
  \par\medskip\noindent
  \fcolorbox{rulegray}{band}{\parbox{\dimexpr\linewidth-2\fboxsep-2\fboxrule}{\vspace{2pt}#1\vspace{2pt}}}%
  \par\medskip}

\newcommand{\vect}[1]{\boldsymbol{#1}}

\title{}
\author{}
\date{}

\begin{document}

\begin{center}
{\LARGE\bfseries 微分方程作业答案}\\[6pt]
{\large Lecture 16（矩阵指数 $e^{Ax}$ 与 $Y'=AY$）\quad Lecture 17--18（非齐次线性系统）}\\[10pt]
{\small 生成日期：2026-10-19\quad｜\quad 四道题的结论均已用 SymPy 符号验算}\\[2pt]
{\small 作业一截止 10-20 10:08\quad｜\quad 作业二截止 10-22 13:30}
\end{center}

\vspace{4pt}

%=====================================================================
\hdr{作业一\ \ Lecture 16：求 $e^{Ax}$，并解 $Y'=AY$}

两种通用做法，按矩阵长得像哪一个选：

\begin{enumerate}[label=(\arabic*)]
\item \textbf{可对角化（A 有 n 个线性无关特征向量）}

  \[ A=P\Lambda P^{-1},\quad \Lambda=\operatorname{diag}(\lambda_1,\dots,\lambda_n)
     \;\Longrightarrow\;
     e^{Ax}=P\,\operatorname{diag}\!\left(e^{\lambda_1 x},\dots,e^{\lambda_n x}\right)P^{-1},
     \qquad Y=e^{Ax}C \]

\item \textbf{若 $A^2=\mu I$（$\mu$ 为常数）}：把级数 $e^{Ax}=\sum_{k\ge 0}\frac{(Ax)^k}{k!}$ 按 $k$ 的奇偶拆成两支，得

  \[
  e^{Ax}=
  \begin{cases}
    \cosh(\sqrt{\mu}\,x)\,I+\dfrac{\sinh(\sqrt{\mu}\,x)}{\sqrt{\mu}}\,A, & \mu>0,\\[10pt]
    \cos(\sqrt{|\mu|}\,x)\,I+\dfrac{\sin(\sqrt{|\mu|}\,x)}{\sqrt{|\mu|}}\,A, & \mu<0.
  \end{cases}
  \]

  （$\mu=0$ 时为 $I+xA$。这条捷径对作业一的第 1、2 题最快，也最不容易算错。）
\end{enumerate}

%---------------------------------------------------------------------
\subsection*{1.\quad $A=\begin{pmatrix}-6&-4\\ 8&6\end{pmatrix}$}

\[
A=\begin{pmatrix}-6&-4\\ 8&6\end{pmatrix},\qquad
\operatorname{tr}A=0,\quad \det A=(-6)(6)-(-4)(8)=-4
\]
特征方程 $\lambda^2-(\operatorname{tr}A)\lambda+\det A=\lambda^2-4=0$，故 $\lambda_1=2,\ \lambda_2=-2$，
且直接算得 $A^2=4I$。对应特征向量为 $\lambda=2:(1,-2)^{\mathsf T}$，$\lambda=-2:(1,-1)^{\mathsf T}$。

由 $A^2=4I$，$e^{Ax}=\cosh 2x\,I+\dfrac{\sinh 2x}{2}A$：

\[
e^{Ax}
=\begin{pmatrix}\cosh 2x-3\sinh 2x & -2\sinh 2x\\[2pt] 4\sinh 2x & \cosh 2x+3\sinh 2x\end{pmatrix}
=\begin{pmatrix}-e^{2x}+2e^{-2x} & -e^{2x}+e^{-2x}\\[2pt] 2e^{2x}-2e^{-2x} & 2e^{2x}-e^{-2x}\end{pmatrix}
\]

\resultband{%
通解 $Y=e^{Ax}C$ 写成分量：
\[
\begin{aligned}
y_1&=C_1\left(-e^{2x}+2e^{-2x}\right)+C_2\left(-e^{2x}+e^{-2x}\right),\\
y_2&=C_1\left(2e^{2x}-2e^{-2x}\right)+C_2\left(2e^{2x}-e^{-2x}\right).
\end{aligned}
\]
等价的特征向量形式（更好记）：\ $Y=a\,e^{2x}\begin{pmatrix}1\\-2\end{pmatrix}+b\,e^{-2x}\begin{pmatrix}1\\-1\end{pmatrix}$。}

%---------------------------------------------------------------------
\subsection*{2.\quad $A=\begin{pmatrix}-4&4\\ -5&4\end{pmatrix}$}

\[
\operatorname{tr}A=0,\qquad \det A=(-4)(4)-4(-5)=4,\qquad
\lambda^2+4=0\ \Longrightarrow\ \lambda=\pm 2i,\qquad A^2=-4I
\]

由 $A^2=-4I$，$e^{Ax}=\cos 2x\,I+\dfrac{\sin 2x}{2}A$：

\[
e^{Ax}=
\begin{pmatrix}\cos 2x-2\sin 2x & 2\sin 2x\\[2pt] -\dfrac{5}{2}\sin 2x & \cos 2x+2\sin 2x\end{pmatrix}
\]

\resultband{%
通解（由特征向量 $(2,\,2+i)^{\mathsf T}$ 的复解拆实部、虚部，或直接取 $e^{Ax}$ 的两列）：
\[
Y=C_1\begin{pmatrix}2\cos 2x\\[2pt] 2\cos 2x-\sin 2x\end{pmatrix}
 +C_2\begin{pmatrix}2\sin 2x\\[2pt] 2\sin 2x+\cos 2x\end{pmatrix}
\]}

%---------------------------------------------------------------------
\subsection*{3.\quad $A=\begin{pmatrix}-24&0&0\\ 4&-2&0\\ 1&0&2\end{pmatrix}$}

$A$ 是下三角矩阵，特征值就是对角元 $\lambda=-24,\,-2,\,2$，两两不同，故可对角化。
另外，这个系统是三角的：$y_1'=-24y_1$ 只含 $y_1$，$y_2'=-2y_2+4y_1$ 只含 $y_1,y_2$，
$y_3'=2y_3+y_1$ 只含 $y_1,y_3$。所以可以按 $y_1\to y_2,y_3$ 的顺序逐层积分（级联），
不必先算特征向量：

\[
\begin{aligned}
y_1'&=-24\,y_1 &&\Longrightarrow\quad y_1=C_1e^{-24x},\\[2pt]
\left(y_2e^{2x}\right)'&=4y_1e^{2x}=4C_1e^{-22x}
 &&\Longrightarrow\quad y_2=C_2e^{-2x}-\tfrac{2}{11}C_1e^{-24x},\\[2pt]
\left(y_3e^{-2x}\right)'&=y_1e^{-2x}=C_1e^{-26x}
 &&\Longrightarrow\quad y_3=C_3e^{2x}-\tfrac{1}{26}C_1e^{-24x}.
\end{aligned}
\]

\resultband{%
指数矩阵（满足 $E(0)=I$ 与 $E'=AE$）：
\[
e^{Ax}=
\begin{pmatrix}
e^{-24x} & 0 & 0\\[4pt]
\dfrac{2}{11}\left(e^{-2x}-e^{-24x}\right) & e^{-2x} & 0\\[6pt]
\dfrac{1}{26}\left(e^{2x}-e^{-24x}\right) & 0 & e^{2x}
\end{pmatrix}
\]
通解 $Y=e^{Ax}C$：
\[
\begin{aligned}
y_1&=C_1e^{-24x},\\
y_2&=\tfrac{2}{11}C_1\left(e^{-2x}-e^{-24x}\right)+C_2e^{-2x},\\
y_3&=\tfrac{1}{26}C_1\left(e^{2x}-e^{-24x}\right)+C_3e^{2x}.
\end{aligned}
\]}

%=====================================================================
\newpage
\hdr{作业二\ \ Lecture 17--18：用常数变易法与待定系数法解非齐次系统}

\[
A=\begin{pmatrix}-2&-2&2\\ -3&1&-4\\ -3&2&-5\end{pmatrix},
\qquad
f(x)=\begin{pmatrix}2x+3\\ 3x-1\\ 3x-2\end{pmatrix},
\qquad
Y=\begin{pmatrix}y_1\\ y_2\\ y_3\end{pmatrix},\qquad Y'=AY+f(x)
\]

\subsection*{2.1\quad 齐次部分}

\[
\det(A-\lambda I)=-(\lambda+1)(\lambda+2)(\lambda+3)\ \Longrightarrow\ \lambda=-1,\,-2,\,-3
\]
对应特征向量
\[
\vect{v}_1=\begin{pmatrix}-2\\5\\4\end{pmatrix},\qquad
\vect{v}_2=\begin{pmatrix}-1\\3\\3\end{pmatrix},\qquad
\vect{v}_3=\begin{pmatrix}0\\1\\1\end{pmatrix}
\]
\[
P=(\vect{v}_1\ \vect{v}_2\ \vect{v}_3)=
\begin{pmatrix}-2&-1&0\\ 5&3&1\\ 4&3&1\end{pmatrix},
\qquad \det P=1,\qquad
P^{-1}=\begin{pmatrix}0&1&-1\\ -1&-2&2\\ 3&2&-1\end{pmatrix}
\]
齐次通解
\[
Y_h=C_1e^{-x}\begin{pmatrix}-2\\5\\4\end{pmatrix}
   +C_2e^{-2x}\begin{pmatrix}-1\\3\\3\end{pmatrix}
   +C_3e^{-3x}\begin{pmatrix}0\\1\\1\end{pmatrix}.
\]

\subsection*{2.2\quad 方法一：待定系数法}

右端是一次多项式，试探 $Y_p=\vect{a}x+\vect{b}$（$\vect a,\vect b$ 为常向量）。代入 $Y'=AY+f(x)$：

\[
\vect a=A(\vect a x+\vect b)+\left(2x+3,\ 3x-1,\ 3x-2\right)^{\mathsf T}
\]
比较 $x$ 的系数与常数项：

\[
A\vect a+\begin{pmatrix}2\\3\\3\end{pmatrix}=\vect 0
\ \Longrightarrow\
A\vect a=\begin{pmatrix}-2\\-3\\-3\end{pmatrix}
\ \Longrightarrow\
\vect a=\begin{pmatrix}1\\0\\0\end{pmatrix},
\]

\[
A\vect b+\begin{pmatrix}3\\-1\\-2\end{pmatrix}=\vect a
\ \Longrightarrow\
A\vect b=\begin{pmatrix}1\\0\\0\end{pmatrix}-\begin{pmatrix}3\\-1\\-2\end{pmatrix}
=\begin{pmatrix}-2\\1\\2\end{pmatrix}
\ \Longrightarrow\
\vect b=\begin{pmatrix}0\\1\\0\end{pmatrix}.
\]

于是 $Y_p=\vect a\,x+\vect b=(x,\ 1,\ 0)^{\mathsf T}$。

同一件事在解耦坐标 $z=P^{-1}Y$ 下更好看：$z_1'+z_1=1$，$z_2'+2z_2=-2x-5$，$z_3'+3z_3=9x+9$，
故 $z_p=\left(1,\ -x-2,\ 3x+2\right)^{\mathsf T}$，再左乘 $P$ 即得 $Y_p=(x,1,0)^{\mathsf T}$。

\subsection*{2.3\quad 方法二：常数变易法}

取基本解矩阵

\[
\Phi(x)=P\,\operatorname{diag}\!\left(e^{-x},e^{-2x},e^{-3x}\right)
=
\begin{pmatrix}-2e^{-x} & -e^{-2x} & 0\\[2pt]
 5e^{-x} & 3e^{-2x} & e^{-3x}\\[2pt]
 4e^{-x} & 3e^{-2x} & e^{-3x}\end{pmatrix},
\qquad
Y=\Phi(x)\left[\vect C+\int \Phi^{-1}(x)f(x)\,dx\right]
\]

先算 $\Phi^{-1}(x)f(x)$：

\[
\Phi^{-1}(x)f(x)=
\begin{pmatrix}e^{x}\\[2pt] (-2x-5)e^{2x}\\[2pt] 9(x+1)e^{3x}\end{pmatrix}
\ \xrightarrow{\ \int dx\ }\
\begin{pmatrix}e^{x}\\[2pt] -(x+2)e^{2x}\\[2pt] (3x+2)e^{3x}\end{pmatrix}
\]

再左乘 $\Phi(x)$：

\[
\Phi(x)\!\int\Phi^{-1}f\,dx
=\begin{pmatrix}-2e^{-x}\cdot e^{x}-e^{-2x}\left[-(x+2)e^{2x}\right]\\[4pt]
 5e^{-x}\cdot e^{x}+3e^{-2x}\left[-(x+2)e^{2x}\right]+e^{-3x}(3x+2)e^{3x}\\[4pt]
 4e^{-x}\cdot e^{x}+3e^{-2x}\left[-(x+2)e^{2x}\right]+e^{-3x}(3x+2)e^{3x}\end{pmatrix}
=\begin{pmatrix}x\\1\\0\end{pmatrix}
\]

与待定系数法的特解完全一致，两法结果相同。

\resultband{%
\textbf{最终答案}（$C_1,C_2,C_3$ 为任意常数）：
\[
\begin{aligned}
y_1&=-2C_1e^{-x}-C_2e^{-2x}+x,\\[2pt]
y_2&=5C_1e^{-x}+3C_2e^{-2x}+C_3e^{-3x}+1,\\[2pt]
y_3&=4C_1e^{-x}+3C_2e^{-2x}+C_3e^{-3x}.
\end{aligned}
\]
即 $Y=Y_h+(x,1,0)^{\mathsf T}$：齐次部分由 $P$ 的列乘 $e^{\lambda x}$ 组成，非齐次部分正是待定系数法得到的 $Y_p$。}

%=====================================================================
\subsection*{3.\quad 验算记录}

四份结论都用 SymPy 做了符号回代，全部通过：

\begin{itemize}
\item 作业一：三个 $e^{Ax}$ 均满足 $E(0)=I$、$E'=AE$、$EA=AE$，且与 \texttt{sympy.exp(A*x)} 逐项化简为同一矩阵；
\item 作业一第 3 题：$A$ 的特征值确为 $\{-24,-2,2\}$；
\item 作业二：通解回代满足 $Y'=AY+f$；特解 $(x,1,0)^{\mathsf T}$ 单独回代同样成立；
\item 作业二：$\Phi(x)$ 的求逆、$\Phi^{-1}f$、积分与 $\Phi\!\int\Phi^{-1}f\,dx$ 四步逐项核对，结果回到 $(x,1,0)^{\mathsf T}$。
\end{itemize}

\end{document}
